\chapter{Допуск off-diagonal linking bridge для карты c0}
\label{ch:v8-baryon-c0-linking-algebra-offdiagonal-bridge-admission-gate}

Direct-sum trace оставил свободное отношение центральных весов. Проверим
минимальную неразделённую архитектуру: исходная и auxiliary прямые становятся
двумя углами одного полного матричного блока
\begin{equation}
 \mathcal L_{min}=M_2(\mathbb C)=
 \begin{pmatrix}
 \mathbb C&\mathbb C\\
 \mathbb C&\mathbb C
 \end{pmatrix}.
 \label{eq:v8-baryon-c0-minimal-linking-algebra}
\end{equation}

\section{Матричная единица как мост}

Положим
\begin{equation}
 P_s=\begin{pmatrix}1&0\\0&0\end{pmatrix},
 \qquad
 P_a=\begin{pmatrix}0&0\\0&1\end{pmatrix},
 \qquad
 E=\begin{pmatrix}0&1\\0&0\end{pmatrix}.
 \label{eq:v8-baryon-c0-linking-matrix-units}
\end{equation}
Тогда
\begin{equation}
 EE^*=P_s,
 \qquad E^*E=P_a,
 \qquad P_sE=E=EP_a.
 \label{eq:v8-baryon-c0-linking-imprimitivity-relations}
\end{equation}
Наличие $E$ связывает два угла: порождённая $*$-алгебра и её центр равны
\begin{equation}
 C^*(P_s,P_a,E)=M_2(\mathbb C),
 \qquad Z(M_2(\mathbb C))=\mathbb C I_2.
 \label{eq:v8-baryon-c0-linking-generated-simple-block}
\end{equation}
Поэтому независимых центральных весов больше нет.

\section{Единый след и нормировка карты}

Единственный нормированный trace полного блока есть
\begin{equation}
 \tau_2(X)=\frac12\Tr_2X,
 \qquad
 \tau_2(P_s)=\tau_2(P_a)=\frac12.
 \label{eq:v8-baryon-c0-linking-unique-trace}
\end{equation}
Для единичных генераторов двух углов индуцированные метрики совпадают:
\begin{equation}
 G_s=\tau_2(P_s)=\frac12,
 \qquad
 G_a=\tau_2(P_a)=\frac12.
 \label{eq:v8-baryon-c0-linking-corner-metrics}
\end{equation}
Условие pullback-изометрии теперь даёт
\begin{equation}
 G_a\kappa^2=G_s
 \quad\Longrightarrow\quad
 \kappa=1
 \qquad(\kappa>0).
 \label{eq:v8-baryon-c0-linking-kappa-one}
\end{equation}

Даже возможная амплитуда моста не возвращает непрерывную свободу, если мост
должен быть imprimitivity matrix unit. Для $E_h=hE$ имеем
\begin{equation}
 E_hE_h^*=h^2P_s,
 \qquad E_h^*E_h=h^2P_a,
 \label{eq:v8-baryon-c0-scaled-bridge-products}
\end{equation}
и условия $E_hE_h^*=P_s$, $E_h^*E_h=P_a$ вместе с $h>0$ дают
\begin{equation}
 h^2=1\quad\Longrightarrow\quad h=1.
 \label{eq:v8-baryon-c0-imprimitivity-fixes-bridge-scale}
\end{equation}

Для сильнейшего внутреннего кандидата $r_*=4$ условный итог архитектуры
равен
\begin{equation}
 c_0=\kappa r_*=4.
 \label{eq:v8-baryon-c0-linking-conditional-value}
\end{equation}

\begin{theorem}[Минимальный linking pass]
Ненулевая imprimitivity-стрелка между двумя одномерными углами порождает
простой $M_2$-блок, устраняет симплекс центральных следовых весов и при
положительной ориентации фиксирует $\kappa=1$. Поэтому минимальная linking
архитектура условно выбирает $c_0=r_*$.
\end{theorem}

\section{Граница происхождения}

Получен допуск архитектуры, но не её наследование. Текущий 42-мерный
центрированный кадр является пространством шумовых и полевых направлений;
в нём ещё не классифицирована gauge-, Real- и grading-совместимая матричная
единица, чей начальный угол есть прямая $r_*$, а конечный --- новый
auxiliary carrier шеститочечного ядра.

\begin{nogocriterion}[Linking block не возникает из обозначения]
Нельзя объявлять два скалярных коэффициента углами одного $M_2$ только ради
получения равных trace weights. Необходимо предъявить внутри объявленного
операторного носителя $E$ с точными endpoint, Real, grading и gauge
условиями и проверить imprimitivity identities.
\end{nogocriterion}

\begin{protocol}\begin{enumerate}
\item минимальная linking-алгебра: \textbf{$M_2(\mathbb C)$};
\item центр: \textbf{$\mathbb C I_2$};
\item нормированный trace: \textbf{единственный};
\item $\kappa$: \textbf{$1$ условно};
\item $c_0$ при $r_*=4$: \textbf{$4$ условно};
\item существующий 42-carrier содержит требуемый $E$: \textbf{не проверено};
\item следующий гейт: \textbf{классификация моста в текущем носителе}.
\end{enumerate}\end{protocol}

Машинный сертификат сохранён в
\path{s2t_v8_baryon_c0_linking_algebra_offdiagonal_bridge_admission_gate_results.json}.